\chapter{Conclusion and perspectives}
\label{chap:conclusion}

\section{Summary}
\label{sec:conc-summary}

We formalised the identification of radiological events as sliding-window
multi-class classification of a noisy gamma-dose signal, and validated a
deep-learning approach to it on a controlled synthetic benchmark. The benchmark
rests on a deliberately simplified signal model --- white Gaussian noise on an
Ornstein--Uhlenbeck baseline, with parameters calibrated to four months of real
2015 recordings and six injected anomaly morphologies. On held-out synthetic
data, all three architectures reach a strong event-level performance (macro $F_1$
from $0.93$ to $1.00$ on a $59$-event segment), with no false alarms. The
residual network is the best overall: a perfect event-level score, the highest
window accuracy, and --- decisively --- the only model besides the attention
hybrid that resists background-noise growth, while the baseline 1D-CNN collapses
to a single class. Self-attention did not improve on the ResNet at the
$512$-sample window scale.

\section{Architecture recommendation}
\label{sec:conc-reco}

For this task we recommend the \textbf{ResNet 1D}: it combines the best accuracy
and robustness with a moderate cost ($\sim 75$k parameters, $\sim 10$ min to
train on a laptop GPU). The 1D-CNN is attractive only where compute is extremely
constrained \emph{and} noise conditions are stable; the CNN+Attention is the
natural candidate if the window is later extended to capture long-range context.

\section{Limitations}
\label{sec:conc-limits}

The results are, by construction, \emph{synthetic}. The signal model is simple in
two ways that matter: its high-frequency noise is independent in time (real
sensor noise is correlated), and its anomaly density is far higher than a real
beacon's. As stressed throughout, this is appropriate for validating the
approach but means the numbers above should not be read as real-world detection
rates. There is also no labelled real test set, so real performance cannot yet be
measured.

\section{Perspectives}
\label{sec:conc-future}

The natural next step --- in the supervisor's words, ``to go further'' --- is to
close the gap to real signals and only then evaluate on them:
\begin{itemize}[leftmargin=*, itemsep=0.25em]
    \item \textbf{More realistic signal models.} Replace the white noise by a
    coloured (e.g.\ autoregressive) process matching the empirical
    autocorrelation of the real residual, lower the anomaly density towards the
    real one, and calibrate amplitudes/durations and the morphology set to events
    actually observed. The very same validation pipeline can then be re-run on
    this harder benchmark.
    \item \textbf{A labelled real test set.} Expert labelling of the real
    recordings would, for the first time, allow event-level precision and recall
    to be measured on real data --- the only sound way to evaluate networks
    trained on a synthetic model.
    \item \textbf{Explainability.} The convolutional models expose their last
    feature map for Grad-CAM \cite{selvaraju2017gradcam}, and the encoder's
    attention maps are directly interpretable; both would let an operator verify
    that a prediction attends to the right part of an event.
    \item \textbf{Longer windows and richer inputs.} Longer windows would give
    the attention model the context it lacked here, and the convolutions already
    accept extra channels, so the companion meteorological signals (temperature,
    humidity, pressure) could be added as inputs.
\end{itemize}

\section{Final word}
\label{sec:conc-final}

On a clean, controlled benchmark the approach works well and the architectural
comparison is informative: depth and group normalisation, not attention, are what
buy robustness at this window length. The honest limitation --- that this is
synthetic validation, not real-world evaluation --- is also the clearest signpost
for the next iteration.
